\documentclass[conference]{IEEEtran}
\IEEEoverridecommandlockouts
\usepackage{cite}
\usepackage{amsmath,amssymb,amsfonts}
\usepackage{algorithmic}
\usepackage{graphicx}
\usepackage{textcomp}
\usepackage{xcolor}
\usepackage[hyphens]{url}
\def\BibTeX{{\rm B\kern-.05em{\sc i\kern-.025em b}\kern-.08em
    T\kern-.1667em\lower.7ex\hbox{E}\kern-.125emX}}
\begin{document}

\title{Efficient Polling-Based Learning Rate Optimization for Neural Networks}

\author{
\IEEEauthorblockN{Luiz Henrique}
\IEEEauthorblockA{\textit{} \\
\textit{Universidade Federal de Pernambuco} \\
Recife, Brazil \\
lhbas@cin.ufpe.br}
\and
\IEEEauthorblockN{Jos\'{e} Ronaldo}
\IEEEauthorblockA{\textit{} \\
\textit{Universidade Federal de Pernambuco} \\
Recife, Brazil \\
jrss@cin.ufpe.br}
}

\maketitle

\begin{abstract}
Selecting an appropriate learning rate is one of the most critical decisions in training neural networks effectively. This work has two objectives: (1) to replicate the Polling Method proposed by Tan et al. on the CIFAR-10 benchmark, adapting it from the original ANN setting to a convolutional neural network trained with Adam, and (2) to propose and evaluate two novel extensions --- a Loss Responsive Learning Rate Scheduler and an Improved Polling method --- against a fixed-rate Adam baseline over 150 epochs. The replicated Polling method matches the baseline at 42.35\% test accuracy without any explicit LR tuning, confirming its effectiveness as an adaptive strategy. The Loss Responsive LR Scheduler surpasses both, achieving 51.06\% best validation accuracy before eventually diverging due to an excessively large learning rate. The Improved Polling method underperforms due to a learning rate collapse caused by selection bias toward smaller gradient steps. Our results reveal fundamental trade-offs between exploration-driven and formula-driven adaptive strategies and highlight open challenges in combining both paradigms.
\end{abstract}

\begin{IEEEkeywords}
learning rate optimization, neural networks, polling method, CIFAR-10, adaptive learning rate, Adam optimizer
\end{IEEEkeywords}

\section{Introduction}

The learning rate (LR) is arguably the most influential hyperparameter in gradient-based training of neural networks. A rate that is too large causes divergence or oscillation near optima, while a rate that is too small leads to slow convergence or entrapment in poor local minima \cite{b1}. Fixed learning rates require expert knowledge or costly grid search, motivating the development of adaptive strategies.

Existing approaches range from formula-based schedulers (step decay, cosine annealing, ReduceLROnPlateau) to optimizers with built-in per-parameter adaptation such as Adam \cite{b1} and RMSProp. A fundamentally different paradigm is \emph{polling optimization}, introduced in \cite{b2}: rather than computing a single LR-scaled update, multiple candidate learning rates are tested at each gradient step and the one that most improves performance on the current batch is selected. This exploration-based approach requires no hand-designed formula and adapts implicitly to the loss landscape.

This work pursues two goals. First, we replicate the Polling Method on CIFAR-10, adapting it from the original ANN and MNIST setting to a convolutional neural network trained with Adam --- a setup that differs from the original in architecture, optimizer, and candidate LR range. Second, we propose two novel extensions (a Loss Responsive LR Scheduler and an Improved Polling method) and compare all four strategies against a fixed-rate Adam baseline. We report training dynamics, convergence speed, and generalization accuracy across 150 epochs.

\section{Proposed Solution}

\subsection{Replicated Method: Polling}

The Polling Method, proposed by Tan et al.~\cite{b2}, replaces the standard single-LR gradient step with an ensemble selection procedure. At every training batch, after computing gradients, $N$ candidate learning rates are applied independently and the model weights that achieve the highest batch accuracy are retained.

Let $\mathcal{C} = \{lr_1, \ldots, lr_N\}$ be the candidate set. At batch $t$, given weights $\theta_t$ and gradients $g_t$:
\begin{equation}
\hat{\theta}_{t,k} = \theta_t - lr_k \cdot g_t, \quad k = 1,\ldots,N
\end{equation}
\begin{equation}
k^{*} = \arg\max_{k} \; \mathrm{acc}(\hat{\theta}_{t,k},\, \mathcal{B}_t)
\end{equation}
\begin{equation}
\theta_{t+1} = \hat{\theta}_{t,k^{*}}
\end{equation}

where $\mathcal{B}_t$ is the current mini-batch. We use $\mathcal{C} = \{10^{-7},\, 10^{-6},\, 10^{-5},\, 10^{-4},\, 10^{-3}\}$ centered around the base LR $= 10^{-5}$. Note that Tan et al.\ originally used LRs in the range $[0.01, 0.05]$ with a simple ANN and SGD-like updates; our candidate set is adapted to the Adam optimizer and CNN setting.

\subsection{Proposed Extensions}

The following two methods are novel contributions of this work, extending the polling paradigm with formula-driven adaptation.

\subsubsection{Loss Responsive LR Scheduler}

The Loss Responsive LR Scheduler adjusts the learning rate based on how much the loss improves over a fixed patience window. An exponential moving average (EMA) smooths the loss signal to reduce stochasticity:

\begin{equation}
\ell_t^{\mathrm{ema}} = \beta \cdot \ell_{t-1}^{\mathrm{ema}} + (1-\beta) \cdot \ell_t, \quad \beta = 0.9
\end{equation}

Every \texttt{patience} steps the improvement ratio is computed as:
\begin{equation}
r = \frac{\ell_{\mathrm{ref}}}{\ell_t^{\mathrm{ema}}}
\end{equation}

Two logarithmic factors are defined, both evaluating to $1$ when $r = 1$:
\begin{align}
f_{\uparrow}(r) &= \log_2(1 + r) \\
f_{\downarrow}(r) &= 2 - \log_2(1 + r)
\end{align}

The choice of $\log_2(1+r)$ is motivated by three design properties: (i)~both factors evaluate to $1$ when $r = 1$, ensuring the learning rate is unaffected when the loss is stagnant; (ii)~$f_{\uparrow}(r) \geq 0$ for all $r \geq 0$, guaranteeing the learning rate remains non-negative regardless of the improvement ratio; and (iii)~the sublinear (logarithmic) growth rate limits how aggressively large improvement ratios drive LR changes, providing more conservative updates than linear or exponential alternatives. Fig.~\ref{fig:lrfactors} illustrates both curves.

\begin{figure}[htbp]
\centerline{\includegraphics[width=\columnwidth]{images/learning_rate_factor_formula.png}}
\caption{LR update factors $f_{\uparrow}(r) = \log_2(1+r)$ (black) and $f_{\downarrow}(r) = 2 - \log_2(1+r)$ (red) as functions of the improvement ratio $r$. Both curves cross at $(1, 1)$: when there is no loss change ($r=1$) neither factor alters the learning rate. $f_{\uparrow}$ starts at the origin $(0, 0)$, is always non-negative, and grows sub-linearly, limiting the magnitude of aggressive LR increases. $f_{\downarrow}$ starts at $(0, 2)$ and decreases symmetrically, reaching zero at $r = 3$.}
\label{fig:lrfactors}
\end{figure}

A direction flag $d \in \{\uparrow, \downarrow\}$ controls whether the LR should increase or decrease. The formula applied at each check is chosen so that $d$ always reflects the true direction of the update:
\begin{equation}
lr \leftarrow
\begin{cases}
lr \cdot f_{\uparrow}(r) & \text{if } d = \uparrow \text{ and } r > 1 \\
lr \cdot f_{\downarrow}(r) & \text{if } d = \downarrow \text{ and } r > 1 \\
lr \cdot f_{\downarrow}(r) & \text{if } d = \uparrow \text{ and } r \leq 1 \\
lr \cdot f_{\uparrow}(r)  & \text{if } d = \downarrow \text{ and } r \leq 1
\end{cases}
\end{equation}

For $r > 1$, $f_{\uparrow}(r) > 1$ and $f_{\downarrow}(r) < 1$; the natural assignment applies and $d$ is unchanged. For $r \leq 1$, $f_{\downarrow}(r) \geq 1$ and $f_{\uparrow}(r) \leq 1$; the formulas are swapped so that $d = \uparrow$ still produces an LR increase and $d = \downarrow$ still produces a decrease. After every non-improving check $d$ alternates to enable oscillatory exploration.

Parameters used: patience $= 30$, $\beta = 0.9$, no LR bounds.

\subsubsection{Improved Polling Method}

The Improved Polling Method combines formula-driven adaptation with candidate selection. Instead of $N$ fixed candidates, it derives two adaptive factors from the current loss improvement ratio and tests three candidates:

\begin{equation}
r = \frac{\ell_{\mathrm{window\_start}}}{\ell_{\mathrm{current}}}
\end{equation}
\begin{align}
f_{\mathrm{up}} &= \log_2(1 + r) \\
f_{\mathrm{down}} &= 2 - \log_2(1 + r)
\end{align}

The same $\log_2(1+r)$ family used in the Loss Responsive Scheduler (Section~II-B1) is applied here, sharing its key properties: both factors equal $1$ at $r = 1$, $f_{\mathrm{up}} \geq 0$ for all $r \geq 0$, and logarithmic growth limits abrupt LR jumps. The three candidates $\{lr \cdot f_{\mathrm{up}},\; lr,\; lr \cdot f_{\mathrm{down}}\}$ are evaluated on the batch, and the one minimizing batch loss is selected. Parameters: patience $= 5$, min\_lr $= 10^{-9}$.

\section{Experiments}

\subsection{Experimental Setup}

All experiments use CIFAR-10 \cite{b3}, consisting of 50{,}000 training and 10{,}000 test images across 10 classes ($32\!\times\!32$ RGB). The full implementation is publicly available in the project repository.\footnote{\url{https://github.com/luiz-linkezio/Efficient-Polling-Based-Learning-Rate-Optimization-for-Neural-Networks}} We reserve 10\% of training data for validation, yielding 45{,}000 training and 5{,}000 validation images.

The model, \texttt{SimpleCIFAR10CNN}, is a lightweight CNN with the following structure: two convolutional blocks (64 and 128 channels, $3\!\times\!3$ kernels, ReLU activations, MaxPool downsampling), a third block of 256 channels followed by AdaptiveAvgPool$(1\!\times\!1)$, and a Linear$(256, 10)$ classification head.

Images are normalized per channel using training set statistics (RGB means: $[125.31,\, 122.95,\, 113.87]$; stds: $[62.99,\, 62.09,\, 66.70]$). All methods use Adam \cite{b1} as the base optimizer with weight decay $5\!\times\!10^{-4}$, batch size 2048, and 150 epochs.

\subsection{Baseline Replication}

The baseline applies Adam with a fixed learning rate of $10^{-5}$ throughout all 150 epochs. This configuration serves as the reference point for all comparisons. As shown in Fig.~\ref{fig:losses}, validation loss decreases monotonically and accuracy improves steadily with no signs of divergence or overfitting, achieving 42.85\% best validation accuracy and 42.91\% test accuracy. The conservative LR ensures stable, predictable convergence, providing a clean lower-bound reference.

\subsection{Polling Method Replication}

The primary experimental goal of this work is to replicate the Polling Method of Tan et al.~\cite{b2} in a new setting. The original paper evaluated it on MNIST with a two-layer ANN and candidate LRs in $[0.01, 0.05]$. Our replication targets CIFAR-10, a harder 10-class benchmark, using a CNN (\texttt{SimpleCIFAR10CNN}) and the Adam optimizer with candidate LRs spanning five orders of magnitude ($10^{-7}$ to $10^{-3}$). Despite these differences in architecture, optimizer, and dataset, the core algorithm is preserved exactly: at each batch, $N$ independent weight updates are evaluated and the one with highest batch accuracy is kept. The replication result (42.35\% test accuracy, closely tracking the baseline of 42.91\%) confirms that the polling selection mechanism is viable and transferable across settings, even when the candidate LR range and optimizer differ substantially from the original.

\section{Results}

Table~\ref{tab:results} summarizes the performance of all four methods. Fig.~\ref{fig:losses} shows training and validation loss curves and Fig.~\ref{fig:lrs} shows the learning rate trajectories across all 150 epochs.

\begin{table}[htbp]
\caption{Performance Comparison Across Methods (150 Epochs)}
\begin{center}
\resizebox{\columnwidth}{!}{%
\begin{tabular}{|l|c|c|c|}
\hline
\textbf{Method} & \textbf{Best Val} & \textbf{Test Acc} & \textbf{Test Loss} \\
\hline
Baseline  & 42.85\%        & 42.91\%          & 1.5585 \\
Polling (base paper) \cite{b2} & 42.36\%    & 42.35\%          & 1.5615 \\
Loss Responsive LR         & \textbf{51.06\%} & \textbf{51.94\%} & \textbf{1.3435} \\
Improved Polling           & 33.38\%        & 34.27\%          & 1.8097 \\
\hline
\end{tabular}%
}
\end{center}
\label{tab:results}
\end{table}

\begin{figure}[htbp]
\centerline{\includegraphics[width=\columnwidth]{images/training_comparison_losses.png}}
\caption{Training and validation loss over 150 epochs for all four methods. The Loss Responsive LR Scheduler reaches the lowest validation loss before diverging, the Polling method closely tracks the Adam baseline, and the Improved Polling method exhibits persistently higher loss due to learning rate collapse.}
\label{fig:losses}
\end{figure}

\begin{figure}[htbp]
\centerline{\includegraphics[width=\columnwidth]{images/training_comparison_LRs.png}}
\caption{Mean selected learning rate per epoch for all four methods. The Loss Responsive LR Scheduler grows aggressively up to ${\approx}\,1.49\!\times\!10^{-1}$ by epoch~150; the Polling method exhibits a natural exploration-to-exploitation decay; the Improved Polling method collapses to near zero by epoch~107; and the baseline maintains a constant LR throughout.}
\label{fig:lrs}
\end{figure}

The Polling Method achieves 42.36\% best validation accuracy and 42.35\% test accuracy, closely matching the baseline. The mean selected LR decreases naturally from ${\approx}\,2.2\!\times\!10^{-4}$ at epoch~1 to ${\approx}\,4\!\times\!10^{-6}$ at epoch~150, exhibiting an implicit exploration-to-exploitation transition without any explicit schedule.

The Loss Responsive LR Scheduler reaches peak validation accuracy of 51.06\% at epoch~55, with the LR at approximately $2\!\times\!10^{-2}$ at that point. The model immediately collapses at epoch~57 (validation accuracy drops to ${\approx}\,9\%$), while the LR continues growing through the remainder of training, reaching a maximum of approximately $1.49\!\times\!10^{-1}$ at epoch~150. The logarithmic update factors keep $f_\uparrow(r) \geq 0$ for all $r \geq 0$, yet the scheduler lacks an explicit upper bound on the LR, so the LR grows unchecked past the accuracy peak until the network diverges to near-random performance.

The Improved Polling method plateaus at 33.38\% validation accuracy. Its LR collapses to near zero by epoch~107 due to a selection bias toward smaller steps at batch level, which reduces local noise but prevents meaningful generalization.

\section{Discussion}

\subsection{Polling as Implicit Adaptation}

The Polling Method demonstrates that a simple batch-level selection over candidate learning rates can match a well-tuned Adam baseline without any prior knowledge of the optimal LR. The natural decay of the mean selected LR from $10^{-4}$ in early epochs to $10^{-5}$ later mirrors the behavior of warm-up schedules, confirming that the loss landscape inherently favors larger exploratory steps in early training and finer steps near convergence.

\subsection{Loss Responsive LR: Effectiveness and Fragility}

The Loss Responsive LR Scheduler achieves the strongest results, benefiting from aggressive LR growth when the model improves rapidly and conservative reduction otherwise. The logarithmic update factors $f_{\uparrow}(r) = \log_2(1+r)$ and $f_{\downarrow}(r) = 2-\log_2(1+r)$ satisfy the three design goals stated in Section~II-B1: both equal $1$ at $r=1$, $f_\uparrow$ is always non-negative for $r \geq 0$, and sublinear growth resists the explosive LR jumps that plagued earlier exponential formulations. Despite this, the model still diverges: the LR grows past the accuracy peak at epoch~55, reaching ${\approx}\,1.49\!\times\!10^{-1}$ by epoch~150, and the network collapses to near-random accuracy at epoch~57. The logarithmic brake slows but does not eliminate the divergence, since the formula remains unbounded above for large improvement ratios. Coupling the LR schedule to validation loss (rather than training loss) or imposing an explicit upper LR bound would be the most direct fix.

\subsection{LR Collapse in Improved Polling}

The Improved Polling method suffers from a systematic collapse: at batch level, smaller LR updates consistently produce lower immediate loss (less gradient noise), creating a feedback loop that drives LR to the minimum bound by epoch~107 regardless of the min\_lr safeguard. This reveals a fundamental tension between batch-level fitness and epoch-level learning: a selection criterion that is locally optimal may be globally harmful. Possible mitigations include periodic LR resets, evaluating candidates over multiple consecutive batches, or imposing a higher-level floor that ignores the batch-loss criterion.

\subsection{Trade-offs}

\begin{itemize}
\item \textbf{Computational cost:} The Polling method requires $N$ forward-backward passes per batch ($5\times$ the baseline). The Loss Responsive LR and Improved Polling add negligible overhead beyond one loss comparison.
\item \textbf{Stability:} The fixed baseline is the most stable. The Loss Responsive LR grows aggressively and still diverges, though the logarithmic factors limit instantaneous jumps and keep $f_\uparrow \geq 0$ for all $r \geq 0$; an explicit upper LR bound remains necessary. Improved Polling is unstable due to LR collapse.
\item \textbf{Generalization:} The Loss Responsive LR achieves the largest gain over the baseline, suggesting it reaches a better-generalizing region of the loss landscape before diverging.
\end{itemize}

\section*{Acknowledgment}

The authors thank the course instructor for guidance and feedback throughout this project.

\begin{thebibliography}{00}
\bibitem{b1} D. P. Kingma and J. Ba, ``Adam: A method for stochastic optimization,'' \textit{arXiv preprint arXiv:1412.6980}, 2014.
\bibitem{b2} C. Y. Tan, W. K. Choong, and C. K. Lau, ``Expediting convergence via polling optimisation for gradient descent in neural networks,'' 2026.
\bibitem{b3} A. Krizhevsky, ``Learning multiple layers of features from tiny images,'' University of Toronto, Tech. Rep., 2009.
\bibitem{b4} L. H. B. A. da Silva and J. R. S. Silva, ``Efficient polling-based learning rate optimization for neural networks,'' GitHub, 2026. [Online]. Available: \url{https://github.com/luiz-linkezio/Efficient-Polling-Based-Learning-Rate-Optimization-for-Neural-Networks}
\end{thebibliography}

\end{document}
